\section{Fundamentos teóricos}

\subsection{Transacciones y propiedades ACID}

El acrónimo ACID se consolida en el trabajo de Haerder y Reuter de
1983~\cite{haerder1983}, que agrupa atomicidad, consistencia, aislamiento y
durabilidad como las propiedades que un gestor transaccional debe garantizar.
El cuadro~\ref{tab:acid} traduce cada una al contexto de PayCore.

\begin{table}[htbp]
\centering
\caption{Propiedades ACID aplicadas a PayCore.}
\label{tab:acid}
\small
\begin{tabularx}{\textwidth}{@{}lXX@{}}
\toprule
\textbf{Propiedad} & \textbf{Qué significa aquí} & \textbf{Mecanismo} \\
\midrule
Atomicidad & El débito y el crédito ocurren los dos o ninguno &
  Una sola transacción SQL con \codigo{COMMIT}/\codigo{ROLLBACK} \\
Consistencia & Ningún saldo negativo; la suma de asientos es cero &
  Restricciones \codigo{CHECK} y un trigger de restricción diferido \\
Aislamiento & Dos transferencias simultáneas no se pisan &
  Bloqueo de fila ordenado, o nivel \codigo{SERIALIZABLE} \\
Durabilidad & Una transferencia confirmada sobrevive a un corte &
  Registro de escritura anticipada (WAL) de PostgreSQL \\
\bottomrule
\end{tabularx}
\end{table}

\subsection{Recuperación por registro de escritura anticipada}

La documentación de PostgreSQL establece que los cambios a los ficheros de
datos deben escribirse solo después de que los registros WAL que los describen
se hayan volcado a almacenamiento permanente, de modo que ante una caída la
base se pueda recuperar a partir del registro~\cite{pg16wal}.

La consecuencia práctica para PayCore es importante y a menudo se malinterpreta:
\textbf{la ventana peligrosa no existe dentro de la base de datos}. Si el
proceso muere antes del \codigo{COMMIT}, PostgreSQL descarta el trabajo no
confirmado. La ventana existe \emph{fuera}, entre el \codigo{COMMIT} y la
respuesta HTTP que el cliente nunca recibe, y se cierra con idempotencia
(sección~\ref{sec:idempotencia}) y no con ACID.

\subsection{Control de concurrencia multiversión}

PostgreSQL implementa internamente tres niveles de aislamiento aunque acepta
los cuatro del estándar SQL: su modo \emph{Read Uncommitted} se comporta como
\emph{Read Committed}, que es la única forma sensata de mapear los niveles
estándar a una arquitectura multiversión~\cite{pg16iso}. El nivel por defecto,
Read Committed, evita lecturas sucias pero \textbf{permite} lecturas no
repetibles, fantasmas y anomalías de serialización. Es decir, el nivel por
defecto no protege por sí solo del escenario del cuadro~\ref{tab:lostupdate}.

La taxonomía de anomalías y la razón por la que los niveles ANSI resultan
insuficientes están analizadas en Berenson \emph{et al.}~\cite{berenson1995},
referencia canónica para explicar por qué el aislamiento por instantáneas no
equivale a la serializabilidad.

\subsection{Aislamiento serializable por instantáneas}

Ports y Grittner describen la implementación de \emph{serializable snapshot
isolation} en PostgreSQL~9.1~\cite{ports2012}. Su resumen reporta un coste de
rendimiento inferior al 7\,\% respecto al aislamiento por instantáneas, y en
la evaluación del microbanco SIBENCH matizan que el seguimiento de
dependencias de lectura supone una sobrecarga de CPU del 10 al 20\,\%.

Conviene subrayar que esas cifras son del artículo y se refieren a sus cargas,
no a PayCore. Las mediciones propias, que arrojan diferencias mucho mayores
bajo contención máxima, están en la sección~\ref{sec:rendimiento}, y la razón
de la discrepancia se discute en la sección~\ref{sec:resultados}.

\subsection{Representación del dinero}

El punto flotante binario definido por IEEE~754 no representa exactamente
valores como $0{,}1$ o $0{,}01$, lo que produce errores de redondeo
acumulativos; Goldberg desarrolla la cuestión en detalle~\cite{goldberg1991}.
La propia documentación de PostgreSQL advierte de forma explícita de que los
números en punto flotante no deberían usarse para manejar dinero debido al
potencial de errores de redondeo~\cite{pg16money}.

Las dos representaciones correctas son enteros en la unidad mínima (céntimos)
o un tipo decimal de precisión arbitraria. PayCore elige la primera, y la
justificación está en la sección~\ref{sec:basededatos}.

\subsection{Contabilidad por partida doble}

La partida doble es la práctica contable estándar desde Pacioli (1494). Fowler
la describe como patrón de software: las cuentas son recipientes cuyo saldo se
deriva de los asientos acumulados, y una transacción contable enlaza
explícitamente los asientos entre cuentas~\cite{fowler2006}.

No es un adorno académico. Square documenta haber construido \emph{Books}, un
servicio de contabilidad por partida doble inmutable, precisamente porque sus
servicios de liquidación mostraban ocasionalmente inconsistencias de datos que
derivaban en consultas de clientes y depósitos retrasados~\cite{square2019}.

\subsection{Latencia de cola}

Dean y Barroso argumentan que en sistemas de servicio la media es una medida
engañosa y que las colas de la distribución dominan la experiencia
real~\cite{dean2013}. Esa es la razón por la que la
sección~\ref{sec:rendimiento} reporta percentiles p50, p95 y p99 junto con el
máximo observado, y nunca la media sola.
